% Pregunta 668 (2008-julio-OEP2008, n.º 20): modelo q con costes de ajuste (Romer). Evolución de q en el tiempo con subvención permanente y transitoria (mismo modelo lineal que 668.tex). Autor: Víctor Gutiérrez Marcos
\documentclass[tikz,border=6pt]{standalone}
\input{tikz-tcee}
\begin{document}
\begin{tikzpicture}[tcee,x=1cm,y=1cm]
  \ejes{$t$}{$q$}{6.8}{5.0}
  \draw[guia] (0,4.4) -- (6.3,4.4);
  \draw[guia,tceerojo] (0,1.4) -- (6.3,1.4);
  \node[etiqueta,anchor=east] at (0,4.4) {$1$};
  \node[etiqueta,anchor=east,text=tceerojo] at (0,1.4) {$1-s$};
  % Antes de la subvención
  \draw[curva1] (0.05,4.4) -- (0.6,4.4);
  \draw[guia] (0.6,4.4) -- (0.6,0) node[below] {$t_0$};
  \draw[guia] (1.68,3.306) -- (1.68,0) node[below] {$T$};
  % Permanente
  \draw[curvanueva,dotted,thick] (0.6,4.4) -- (0.6,2.9);
  \draw[curvanueva] plot[smooth] coordinates {(0.600,2.900) (0.825,2.724) (1.050,2.568) (1.275,2.431) (1.500,2.310) (1.725,2.203) (1.950,2.109) (2.175,2.025) (2.400,1.952) (2.625,1.887) (2.850,1.830) (3.075,1.779) (3.300,1.735) (3.525,1.695) (3.750,1.661) (3.975,1.630) (4.200,1.603) (4.425,1.579) (4.650,1.558) (4.875,1.540) (5.100,1.523) (5.325,1.509) (5.550,1.496) (5.775,1.485) (6.000,1.475)};
  \node[etiqueta,anchor=south,text=tceerojo] at (4.4,1.62) {permanente};
  % Transitoria
  \draw[very thick,dashed,black] plot[smooth] coordinates {(0.600,3.352) (0.735,3.281) (0.870,3.227) (1.005,3.190) (1.140,3.172) (1.275,3.172) (1.410,3.193) (1.545,3.237) (1.680,3.306) (1.907,3.436) (2.135,3.550) (2.362,3.651) (2.589,3.740) (2.817,3.818) (3.044,3.887) (3.272,3.948) (3.499,4.002) (3.726,4.049) (3.954,4.091) (4.181,4.127) (4.408,4.160) (4.636,4.188) (4.863,4.213) (5.091,4.235) (5.318,4.255) (5.545,4.272) (5.773,4.287) (6.000,4.301)};
  \node[etiqueta,anchor=north] at (4.2,3.95) {transitoria};
  \node[punto,fill=tceerojo] at (0.6,2.9) {};
  \node[punto] at (0.6,3.352) {};
  \node[punto] at (1.68,3.306) {};
\end{tikzpicture}
\end{document}
